%
%  chapter_deepdive.tex  —  versão 2 (reescrita estrutural)
%  Aprofundamento Técnico: Por Detrás do Play4Change
%
%  Estrutura: geral → específico, um conceito por secção,
%  organizado pelos três pilares do projeto.
%  Linguagem: português europeu académico.
%
\chapter{Aprofundamento Técnico: Por Detrás do Sistema}
\label{cha:deepdive}

% ═══════════════════════════════════════════════════════════════════
\section{Enquadramento e Estrutura do Capítulo}
\label{sec:enquadramento}
% ═══════════════════════════════════════════════════════════════════

Este capítulo percorre o \textit{Play4Change} de dentro para fora.
Em vez de enumerar funcionalidades, parte da arquitectura geral e aprofunda
cada subsistema por ordem de dependência: primeiro o que todos os fluxos
partilham (autenticação, \textit{pipeline} HTTP), depois o que é específico
de cada pilar do projeto.

A leitura está organizada da seguinte forma:

\begin{enumerate}
  \item \textbf{Arquitectura geral} — os três clientes, os serviços de
    suporte e as fronteiras entre camadas.
  \item \textbf{Autenticação} — o protocolo \textit{magic link},
    armazenamento seguro de \textit{tokens} e renovação silenciosa.
  \item \textbf{Pipeline HTTP do servidor} — as catorze camadas que todo
    o pedido atravessa antes de chegar ao controlador.
  \item \textbf{Os três pilares} — mapeamento dos objetivos pedagógicos
    para subsistemas de engenharia.
  \item \textbf{Pilar I — Envolvimento} — atribuição diária de tarefas,
    pontos e sequências.
  \item \textbf{Pilar II — Personalização (criação com IA)} — extração
    de conteúdo, geração de tarefas com LLM e engenharia de
    \textit{prompts}.
  \item \textbf{Pilar II — Personalização (pesquisa vectorial)} —
    \textit{embeddings}, similaridade do cosseno, índice IVFFlat e as
    três estratégias de reutilização.
  \item \textbf{Pilar II — Personalização (remediação adaptativa)} —
    deteção de dificuldade e ramos de recuperação.
  \item \textbf{Pilar III — Reconhecimento} — emissão de \textit{badges}.
  \item \textbf{Arquitetura do cliente móvel} — Kotlin Multiplatform,
    Decompose, Koin e padrão MVI.
  \item \textbf{Observabilidade} — métricas, alertas e \textit{dashboards}.
  \item \textbf{Segurança} — modelo de ameaças e controlos OWASP.
  \item \textbf{Superfície completa da API REST}.
\end{enumerate}

% ═══════════════════════════════════════════════════════════════════
\section{Arquitectura Geral da Plataforma}
\label{sec:arquitectura_geral}
% ═══════════════════════════════════════════════════════════════════

O \textit{Play4Change} é composto por três camadas principais
(Figura~\ref{fig:arquitectura_geral}): dois clientes com superfícies
distintas, uma API REST centralizada e um conjunto de serviços de suporte.

\begin{figure}[htbp]
  \centering
  \begin{tikzpicture}[>=Stealth, font=\small,
    tier/.style={draw, fill=#1, rounded corners=5pt,
                 minimum width=3.2cm, minimum height=1.1cm, align=center},
    svc/.style={draw, fill=gray!12, rounded corners=3pt,
                minimum width=2.4cm, minimum height=0.8cm, align=center},
    lbl/.style={font=\bfseries\footnotesize, color=gray!60}
  ]
    % Clients
    \node[tier=blue!12]   (web) at (-4, 0) {Portal Web\\{\small React + TypeScript}\\{\footnotesize Administrador}};
    \node[tier=orange!15] (mob) at ( 4, 0) {App Móvel\\{\small KMP + Compose}\\{\footnotesize Aprendente}};

    % API
    \node[tier=green!14]  (api) at (0,-2.5){API REST\\{\small Spring Boot}\\{\footnotesize HTTPS / JSON / SSE}};

    % Services
    \node[svc] (pg)    at (-4.5,-5){PostgreSQL\\{\tiny + pgvector}};
    \node[svc] (minio) at (-1.5,-5){MinIO\\{\tiny S3}};
    \node[svc] (ai)    at ( 1.5,-5){Mistral AI\\{\tiny LangChain4j}};
    \node[svc] (resend)at ( 4.5,-5){Resend\\{\tiny E-mail}};

    % Labels
    \node[lbl] at (-6.2, 0)   {Clientes};
    \node[lbl] at (-6.2,-2.5) {Servidor};
    \node[lbl] at (-6.2,-5)   {Serviços};

    % Edges
    \draw[->] (web) -- node[left, font=\tiny]{HTTPS + JWT} (api);
    \draw[->] (mob) -- node[right,font=\tiny]{HTTPS + JWT} (api);
    \draw[->] (api) -- (pg);
    \draw[->] (api) -- (minio);
    \draw[->] (api) -- (ai);
    \draw[->] (api) -- (resend);
    \draw[->, dashed, color=blue!50] (api) to[bend left=12]
      node[above,font=\tiny]{SSE} (web);
  \end{tikzpicture}
  \caption{Arquitectura de três camadas do \textit{Play4Change}}
  \label{fig:arquitectura_geral}
\end{figure}

\paragraph{Portal Web de Administração.}
Construído em React 18 com TypeScript e TanStack Query para gestão do
estado de servidor.
É reservado exclusivamente aos administradores de conteúdo: criação de
tópicos a partir de URL ou PDF, monitorização do progresso de geração em
tempo real via \textit{Server-Sent Events}, edição de tarefas e gestão de
utilizadores.
O administrador nunca interage com a aplicação móvel — são superfícies
completamente separadas com papéis distintos.

\paragraph{Aplicação Móvel.}
Construída com Kotlin Multiplatform e Compose Multiplatform, partilhando
a totalidade do código de negócio, rede e interface entre Android e iOS.
É a superfície primária dos aprendentes: inscrição em tópicos, resolução
de tarefas diárias, progressão no mapa de aprendizagem e remediação
adaptativa quando é detetada dificuldade.

\paragraph{API REST Spring Boot.}
Serve ambos os clientes com o mesmo conjunto de \textit{endpoints},
diferenciados por papel (\texttt{ADMIN} vs.\ utilizador autenticado).
Implementa a Arquitectura Hexagonal~\cite{cockburn2005hexagonal}, isolando
a lógica de domínio de todas as dependências externas através de
interfaces de porta.

\paragraph{Serviços de suporte.}
PostgreSQL (com a extensão \textit{pgvector} para pesquisa vectorial)
persiste toda a informação transacional e os \textit{embeddings} das
tarefas.
MinIO armazena os conteúdos de fonte dos tópicos (texto extraído e PDF
original).
A API Mistral fornece tanto o modelo de linguagem para geração de tarefas
como o modelo de \textit{embedding} para representação semântica.
O Resend entrega os \textit{e-mails} de \textit{magic link}.

% ═══════════════════════════════════════════════════════════════════
\section{Fundação Transversal: Autenticação Sem Palavra-Passe}
\label{sec:autenticacao}
% ═══════════════════════════════════════════════════════════════════

Antes de qualquer funcionalidade, o utilizador — seja administrador no
portal web ou aprendente na aplicação móvel — tem de se autenticar.
O \textit{Play4Change} elimina completamente as palavras-passe, adoptando
um protocolo de \textit{magic link}~\cite{owasp2021a07}: o utilizador
fornece apenas o seu endereço de correio eletrónico e recebe uma ligação
de autenticação de utilização única.

\subsection*{Porque não palavras-passe}

As palavras-passe introduzem três superfícies de ataque que o
\textit{magic link} elimina por construção: armazenamento de hashes
suscetíveis a ataques de dicionário, reutilização de credenciais entre
serviços e \textit{phishing} de credenciais.
A ligação de autenticação é válida durante 15 minutos, de utilização
única e armazenada apenas como resumo SHA-256 — mesmo que a base de dados
seja comprometida, o valor armazenado não pode ser usado diretamente.

\subsection*{O protocolo em quatro passos}

\begin{figure}[htbp]
  \centering
  \begin{tikzpicture}[>=Stealth, font=\small,
    act/.style={draw, fill=#1, rounded corners=3pt,
                minimum width=2.6cm, minimum height=0.7cm, align=center}
  ]
    \node[act=blue!10]   (cli) at (0,   0) {Cliente\\(web / móvel)};
    \node[act=green!10]  (api) at (5.5, 0) {Spring Boot API};
    \node[act=orange!10] (ext) at (10.5,0) {Resend / PostgreSQL};

    % Lifelines
    \foreach \x in {0,5.5,10.5} \draw[dashed] (\x,-0.4) -- (\x,-10.2);

    % Step 1
    \draw[->] (0,-1.0) -- node[above,font=\tiny]{(1) POST /auth/magic-link \{email\}} (5.5,-1.0);
    \draw[->] (5.5,-1.7) -- node[above,font=\tiny]{enviar e-mail com ligação} (10.5,-1.7);
    \draw[->] (5.5,-2.4) -- node[above,font=\tiny]{guardar SHA-256 TTL 15 min} (10.5,-2.4);
    \draw[->] (5.5,-3.1) -- node[above,font=\tiny]{(1R) 202 Accepted} (0,-3.1);

    % Step 2
    \draw[->] (0,-4.2) -- node[above,font=\tiny]{(2) GET /auth/verify?token=X} (5.5,-4.2);
    \draw[->] (5.5,-4.9) -- node[above,font=\tiny]{(2R) 302 → play4change://auth/verify?token=X} (0,-4.9);

    % Step 3
    \draw[->] (0,-6.2) -- node[above,font=\tiny]{(3) POST /auth/magic-link/verify \{token\}} (5.5,-6.2);
    \draw[->] (5.5,-6.9) -- node[above,font=\tiny]{verificar SHA-256, marcar consumido} (10.5,-6.9);
    \draw[->] (5.5,-7.6) -- node[above,font=\tiny]{(3R) 200 OK \{accessToken, refreshToken\}} (0,-7.6);

    % Step 4
    \draw[->] (0,-8.8) -- node[above,font=\tiny]{(4) pedido protegido + Bearer token} (5.5,-8.8);
    \draw[->] (5.5,-9.5) -- node[above,font=\tiny]{(4R) 200 OK / 401 → renovação automática} (0,-9.5);
  \end{tikzpicture}
  \caption{Protocolo de autenticação \textit{magic link}: quatro passos}
  \label{fig:auth_sequence}
\end{figure}

\paragraph{Passo 1 — Pedido de ligação.}
O cliente submete o endereço de correio eletrónico.
O servidor gera um \textit{token} aleatório criptograficamente seguro,
armazena o seu resumo SHA-256 com TTL de 15 minutos, envia o \textit{e-mail}
via API Resend e responde \texttt{202 Accepted} — \emph{sempre a mesma
resposta}, independentemente de o endereço existir ou não, impedindo
enumeração de utilizadores~\cite{owasp2021a05}.
Um \texttt{RateLimitFilter} impõe um máximo de três pedidos por IP por
hora.

\paragraph{Passo 2 — Clique na ligação.}
O \textit{e-mail} contém \texttt{GET /auth/verify?token=\{raw\}}.
O servidor emite \texttt{302 Found} para o esquema URI personalizado
\texttt{play4change://auth/verify?token=\{raw\}}.
No mobile, o sistema operativo interceta o esquema e abre a aplicação.
No web, a \texttt{AuthCallbackPage} processa o redirecionamento.

\paragraph{Passo 3 — Verificação e emissão de \textit{tokens}.}
O cliente envia o \textit{token} em claro via \texttt{POST /auth/magic-link/verify}.
O servidor recalcula o SHA-256, valida a existência e a não-expiração,
marca-o como consumido e emite um par de \textit{tokens}:
um \textit{token} de acesso JWT (15--30 minutos) e um \textit{token} de
atualização opaco com TTL de 7 dias.

\paragraph{Passo 4 — Pedidos autenticados e renovação silenciosa.}
Todos os pedidos subsequentes transportam o JWT no cabeçalho
\texttt{Authorization: Bearer}.
Quando o JWT expira, o cliente renova-o silenciosamente sem interromper
o utilizador (detalhado abaixo por plataforma).

\subsection*{Armazenamento seguro de \textit{tokens} por plataforma}

Os \textit{tokens} não podem ser armazenados em texto simples; cada
plataforma usa o mecanismo de segurança nativo disponível
(Figura~\ref{fig:token_storage}).

\begin{figure}[htbp]
  \centering
  \begin{tikzpicture}[>=Stealth, font=\small,
    plat/.style={draw, fill=#1, rounded corners=5pt,
                 minimum width=4.5cm, minimum height=3.4cm, align=center,
                 text width=4.2cm},
  ]
    \node[plat=blue!8] (web) at (-3.5,0) {%
      \textbf{Portal Web}\\[4pt]
      \textit{Access token}:\\
      \texttt{sessionStorage}\\
      (limpo ao fechar o separador)\\[4pt]
      \textit{Refresh token}:\\
      \textit{cookie} HTTP\\
      \texttt{SameSite=Strict}\\
      \texttt{Secure}\\[4pt]
      Renovação: interceptor Axios\\
      + fila de pedidos pendentes
    };
    \node[plat=green!10] (and) at ( 1.5,0) {%
      \textbf{Android}\\[4pt]
      \texttt{EncryptedSharedPreferences}\\
      (Jetpack Security)\\[4pt]
      Chaves: AES-256-SIV\\
      Valores: AES-256-GCM\\
      Chave-mestra no\\
      \textit{Android Keystore}\\[4pt]
      Renovação: plugin Ktor\\
      \texttt{Auth \{bearer\{\}\}}
    };
    \node[plat=gray!12] (ios) at ( 6.5,0) {%
      \textbf{iOS}\\[4pt]
      \textit{Keychain} (Security.framework)\\
      via Kotlin/Native CInterop\\[4pt]
      \texttt{kSecUseDataProtection}\\
      \texttt{Keychain = true}\\
      \texttt{kSecAttrAccessible}\\
      \texttt{AfterFirstUnlock}\\[4pt]
      Renovação: plugin Ktor\\
      \texttt{Auth \{bearer\{\}\}}
    };
  \end{tikzpicture}
  \caption{Mecanismos de armazenamento seguro de \textit{tokens} por plataforma}
  \label{fig:token_storage}
\end{figure}

Na aplicação móvel, a interface \texttt{TokenStorage} abstrai ambas as
implementações.
O módulo Koin da plataforma (\texttt{platformModule}) regista a
implementação correta em tempo de compilação — o código partilhado nunca
depende de uma implementação concreta.
Quando o servidor devolve \texttt{401}, o plugin Ktor \texttt{Auth}
invoca automaticamente \texttt{POST /auth/refresh}; em caso de falha
não transitória, \texttt{SessionEventBus.sessionExpired()} redireciona
o utilizador para o ecrã de autenticação sem destruir a \textit{Activity}
Android.

% ═══════════════════════════════════════════════════════════════════
\section{O \textit{Pipeline} de Processamento de Pedidos HTTP}
\label{sec:pipeline_http}
% ═══════════════════════════════════════════════════════════════════

Todo o pedido HTTP — independentemente de origem (web ou móvel) ou
destino (rota pública ou protegida) — percorre a mesma sequência de
catorze camadas no servidor antes de atingir o controlador de domínio.
Compreender esta sequência é essencial para entender onde cada
preocupação transversal (segurança, rastreabilidade, métricas) é tratada.

\begin{figure}[htbp]
  \centering
  \begin{tikzpicture}[>=Stealth, font=\small,
    layer/.style={draw, fill=#1, rounded corners=3pt,
                  minimum width=11cm, minimum height=0.78cm,
                  align=center, text width=10.6cm},
    arr/.style={->, thick, color=gray!55}
  ]
    \node[layer=orange!18]  at (0, 0.0)  {\textbf{1. Nginx}
      — terminação TLS, cabeçalhos de segurança, \texttt{proxy\_buffering off} (SSE)};
    \node[layer=yellow!22]  at (0,-0.95) {\textbf{2. Tomcat} (\textit{embedded})
      — aceitação TCP, gestão do conjunto de fios de execução};
    \node[layer=red!14]     at (0,-1.90) {\textbf{3. \texttt{RateLimitFilter}}
      \small[\texttt{HIGHEST\_PRECEDENCE}]\ — limite de pedidos \texttt{/auth/**} por IP → \texttt{429}};
    \node[layer=red!9]      at (0,-2.85) {\textbf{4. \texttt{RequestIdFilter}}
      \small[\texttt{@Order(1)}]\ — gera \texttt{X-Request-Id}, injeta no MDC (rastreabilidade)};
    \node[layer=blue!10]    at (0,-3.80) {\textbf{5. Spring Security — \texttt{CorsFilter}}
      — valida origem, emite \texttt{Access-Control-*}, responde \texttt{OPTIONS}};
    \node[layer=blue!16]    at (0,-4.75) {\textbf{6. \texttt{JwtAuthFilter}}
      — analisa \texttt{Bearer}, povoa \texttt{SecurityContextHolder} (ou \texttt{401})};
    \node[layer=blue!22]    at (0,-5.70) {\textbf{7. Spring Security — Autorização}
      — \texttt{/admin/**} → \texttt{ADMIN}; resto → \texttt{authenticated()}; falha → \texttt{403}};
    \node[layer=green!12]   at (0,-6.65) {\textbf{8. \texttt{WebMvcMetricsFilter}}
      — regista \texttt{http\_server\_requests\_seconds} no Prometheus (automático)};
    \node[layer=purple!10]  at (0,-7.60) {\textbf{9. \texttt{DispatcherServlet}}
      — mapeamento de controlador, negociação de conteúdo, resolução de argumentos};
    \node[layer=teal!10]    at (0,-8.55) {\textbf{10. \texttt{LoggingInterceptor.preHandle()}}
      — grava \textit{timestamp} de início em \texttt{request.setAttribute}};
    \node[layer=teal!16]    at (0,-9.50) {\textbf{11. Controlador + Bean Validation (\texttt{@Valid})}
      — desserialização Jackson, restrições JSR-380 → \texttt{400} em caso de falha};
    \node[layer=teal!22]    at (0,-10.45){\textbf{12. Serviço de Aplicação}
      — lógica de domínio com \texttt{Either<AppError,T>} (Arrow)};
    \node[layer=teal!16]    at (0,-11.40){\textbf{13. \texttt{LoggingInterceptor.afterCompletion()}}
      — regista método, caminho, estado HTTP e duração em ms};
    \node[layer=orange!12]  at (0,-12.35){\textbf{14. \texttt{@RestControllerAdvice}}
      — mapeia exceções não capturadas para respostas de erro JSON estruturadas};

    % downward arrows on the right margin
    \foreach \y in {-0.48,-1.43,-2.38,-3.33,-4.28,-5.23,-6.18,-7.13,
                    -8.08,-9.03,-9.98,-10.93,-11.88} {
      \draw[arr] (5.8,\y) -- (5.8,\y-0.50);
    }
  \end{tikzpicture}
  \caption{As catorze camadas do \textit{pipeline} de pedido-resposta do servidor Spring Boot}
  \label{fig:pipeline_http}
\end{figure}

A lógica de domínio só é alcançada na camada 12.
Todas as preocupações transversais — segurança, rastreabilidade, métricas,
tratamento de erros — estão encapsuladas nas camadas anteriores, sem
qualquer código de infraestrutura nos serviços de aplicação.

\subsection*{Processamento assíncrono paralelo}

Os pedidos de criação de tópico não bloqueiam o fio de execução do
Tomcat enquanto a IA gera as tarefas.
Um \texttt{ThreadPoolTaskExecutor} dedicado (\texttt{generationExecutor}),
com 3 fios de execução e fila de capacidade 25, recebe o trabalho de
geração imediatamente após o registo do tópico na base de dados:

\begin{verbatim}
corePoolSize = 3  |  maxPoolSize = 3  |  queueCapacity = 25
setThreadNamePrefix("gen-")    // visível em JVM thread dumps
\end{verbatim}

Assim, a criação de um tópico devolve \texttt{201 Created} em milissegundos;
o progresso é depois comunicado ao portal web via SSE.
Quando a fila atinge a capacidade, novos pedidos de geração são rejeitados
com \texttt{503} e o alerta Prometheus
\texttt{GenerationThreadPoolQueueSaturated} é disparado.

% ═══════════════════════════════════════════════════════════════════
\section{Os Três Pilares: Mapeamento para Subsistemas de Engenharia}
\label{sec:tres_pilares}
% ═══════════════════════════════════════════════════════════════════

O \textit{Play4Change} está organizado em torno de três pilares
pedagógicos, cada um com um subsistema de engenharia correspondente
(Figura~\ref{fig:pilares_subsistemas}).

\begin{figure}[htbp]
  \centering
  \begin{tikzpicture}[>=Stealth, font=\small,
    pil/.style={draw, fill=#1, rounded corners=5pt,
                minimum width=3.6cm, minimum height=1.2cm, align=center},
    sub/.style={draw, fill=#1, rounded corners=3pt,
                minimum width=3.6cm, minimum height=0.8cm,
                align=center, text width=3.4cm}
  ]
    \node[pil=blue!12]   (eng) at (-4.5,0) {\textbf{Envolvimento}\\gamificação};
    \node[pil=orange!15] (per) at (   0,0) {\textbf{Personalização}\\IA adaptativa};
    \node[pil=green!12]  (rec) at ( 4.5,0) {\textbf{Reconhecimento}\\progresso};

    \node[sub=blue!8]   (e1) at (-4.5,-2.0){\texttt{TaskAssignment}\\1 tarefa/dia/tópico};
    \node[sub=blue!5]   (e2) at (-4.5,-3.2){Pontos + \textit{streak}\\diário};

    \node[sub=orange!10] (p1) at (0,-2.0){LangChain4j + Mistral\\geração de tarefas};
    \node[sub=orange!7]  (p2) at (0,-3.2){pgvector\\similaridade vectorial};
    \node[sub=orange!5]  (p3) at (0,-4.4){\texttt{HandleStruggleService}\\remediação adaptativa};

    \node[sub=green!8]  (r1) at (4.5,-2.0){\texttt{BadgeIssuanceService}\\critérios configuráveis};
    \node[sub=green!5]  (r2) at (4.5,-3.2){\textit{Badges} digitais\\verificáveis (PostgreSQL)};

    \draw[->] (eng) -- (e1); \draw[->] (e1) -- (e2);
    \draw[->] (per) -- (p1); \draw[->] (p1) -- (p2); \draw[->] (p2) -- (p3);
    \draw[->] (rec) -- (r1); \draw[->] (r1) -- (r2);
  \end{tikzpicture}
  \caption{Os três pilares pedagógicos e os subsistemas de engenharia correspondentes}
  \label{fig:pilares_subsistemas}
\end{figure}

As secções seguintes aprofundam cada pilar, começando pelo mais simples
(Envolvimento) e progredindo para o mais complexo (Personalização).

% ═══════════════════════════════════════════════════════════════════
\section{Pilar I — Envolvimento: Atribuição Diária e Gamificação}
\label{sec:pilar_envolvimento}
% ═══════════════════════════════════════════════════════════════════

O pilar de Envolvimento garante que o aprendente tem sempre uma tarefa
disponível e que o progresso acumulado é visível e recompensador.

\subsection*{Inscrição num tópico}

O aprendente inicia a sua jornada no ecrã \textbf{Explorar} da aplicação
móvel, seleciona um tópico ativo e toca em «Inscrever»:

\begin{verbatim}
POST /user/enroll
Authorization: Bearer <token>
{ "topicId": "<uuid>" }
\end{verbatim}

O \texttt{EnrollmentService} aplica quatro guardas em sequência antes de
criar a inscrição: o tópico tem de estar \texttt{ACTIVE}; não pode ter
expirado; todos os pré-requisitos têm de estar concluídos; e não pode
existir já uma inscrição ativa.
Em caso de sucesso, cria imediatamente a primeira \texttt{TaskAssignment}
para \texttt{dayIndex = 0}.
A ordem das opções de cada tarefa é misturada deterministicamente com
\texttt{TaskShuffleSeed.shuffleOptions(optionCount, userId, taskId,
enrollmentId)} — o mesmo utilizador vê sempre a mesma ordenação ao
regressar, mas diferentes utilizadores veem ordenações diferentes,
eliminando o reconhecimento de padrões por posição.

\subsection*{Obtenção da tarefa diária}

A aplicação móvel pede a tarefa do dia ao servidor com:

\begin{verbatim}
GET /tasks/today?topicId={id}
Authorization: Bearer <token>
X-Timezone: Europe/Lisbon
\end{verbatim}

O cabeçalho \texttt{X-Timezone} é fundamental: as atribuições são geradas
por dia de calendário no fuso horário \emph{do aprendente}, não do
servidor.
Sem este cabeçalho, um utilizador em Lisboa às 23h30 poderia receber a
tarefa do dia seguinte se o servidor estivesse num fuso UTC.

O servidor pode devolver três estados especiais além de \texttt{200 OK}:
\texttt{202 Accepted} (geração ainda em curso),
\texttt{404 Not Found} (tarefa já concluída hoje) e
\texttt{409 Conflict} com o cabeçalho
\texttt{X-Open-Struggle-Enrollment} (existe uma sessão de dificuldade
ativa — o cliente é redirecionado para ela em vez de mostrar a tarefa).

\subsection*{Pontos e sequências diárias}

Ao submeter uma resposta correta pela primeira vez num dia, o servidor
atribui \texttt{TaskTemplate.pointsReward} pontos e incrementa
\texttt{Enrollment.streakDays}.
Os pontos são acumulados em \texttt{UserProfile.totalPoints} e o nível
do utilizador é calculado a partir deste total.
Submissões erradas decrementam \texttt{wrongAttemptCount} mas não penalizam
pontos — o sistema é encorajador, não punitivo.

% ═══════════════════════════════════════════════════════════════════
\section{Pilar II — Personalização: Criação de Conteúdo com IA}
\label{sec:pilar_criacao_ia}
% ═══════════════════════════════════════════════════════════════════

A personalização começa na criação do conteúdo.
Em vez de tarefas genéricas, o \textit{Play4Change} gera tarefas
específicas para cada tópico usando um modelo de linguagem de grande
dimensão (LLM), seguindo o enquadramento de engenharia de
\textit{prompts} de quatro componentes da Google~\cite{google2024promptguide}:
\textbf{Persona}, \textbf{Tarefa}, \textbf{Contexto} e \textbf{Formato}.

\subsection*{Criação de tópico pelo administrador}

O administrador acede ao portal web e fornece a fonte de conteúdo:
uma URL pública ou um ficheiro PDF.

\begin{verbatim}
POST /admin/topics
Authorization: Bearer <token_admin>
Content-Type: application/json

{
  "title":      "Mobilidade Urbana Sustentável",
  "url":        "https://example.org/guia-mobilidade",
  "taskCount":  15,
  "difficulty": "BEGINNER",
  "language":   "pt-PT"
}
\end{verbatim}

O servidor responde \texttt{201 Created} imediatamente com o
\texttt{topicId} e o estado \texttt{PENDING}.
A geração decorre integralmente em segundo plano no
\texttt{generationExecutor}.

\subsection*{As fases de geração}

\begin{figure}[htbp]
  \centering
  \begin{tikzpicture}[>=Stealth, font=\small,
    ph/.style={draw, fill=yellow!18, rounded corners=4pt,
               minimum width=2.3cm, minimum height=1.0cm, align=center},
    ok/.style={draw, fill=green!15, rounded corners=4pt,
               minimum width=2.3cm, minimum height=1.0cm, align=center},
    err/.style={draw, fill=red!15,  rounded corners=4pt,
               minimum width=2.3cm, minimum height=1.0cm, align=center}
  ]
    \node[ph]  (ing) at (0,   0) {INGESTION\\{\tiny extração}};
    \node[ph]  (ana) at (3.3, 0) {ANALYSIS\\{\tiny objetivo}};
    \node[ph]  (gen) at (6.6, 0) {GENERATION\\{\tiny LLM}};
    \node[ph]  (idx) at (9.9, 0) {INDEXING\\{\tiny embeddings}};
    \node[ok]  (act) at (9.9,-2.2){ACTIVE};
    \node[err] (fail)at (3.3,-2.2){FAILED};

    \draw[->] (ing) -- (ana);
    \draw[->] (ana) -- (gen);
    \draw[->] (gen) -- (idx);
    \draw[->] (idx) -- (act);
    \draw[->, dashed, red] (gen) -- node[right,font=\tiny]{falha LLM} (fail);
    \draw[->, dashed, red] (ana) -- node[left, font=\tiny]{\textit{timeout}} (fail);
  \end{tikzpicture}
  \caption{Fases do \textit{pipeline} de geração assíncrona de tópicos}
  \label{fig:fases_geracao}
\end{figure}

\paragraph{INGESTION.}
O \texttt{ContentExtractorPort} obtém o conteúdo da URL (remove
marcações HTML) ou extrai texto do PDF (Apache PDFBox).
O texto é enviado para o MinIO (\texttt{topics/\{id\}/content.txt}).
Um evento SSE \texttt{phase-change} é publicado ao portal web.

\paragraph{ANALYSIS.}
Uma análise preliminar extrai o objetivo do módulo — o enunciado de
aprendizagem que o \textit{prompt} de geração usará como âncora
contextual.

\paragraph{GENERATION.}
A chamada central ao LLM.
O \texttt{LangChain4jTaskGenerationAdapter} monta duas mensagens:

\begin{description}
  \item[Mensagem de sistema (Persona + Formato):]
\begin{verbatim}
You are an expert educational content designer
specialising in gamified learning experiences.
OUTPUT FORMAT: valid JSON array only. Schema:
[{ "title":"...", "description":"...", "hint":"...",
   "pointsReward": N, "options":["correta","e1","e2","e3"],
   "correctAnswerIndex": 0 }]
\end{verbatim}
  A resposta correta está \emph{sempre} no índice 0 — o servidor
  embaralha as opções antes de as enviar ao cliente, garantindo que
  o LLM nunca é responsável pela aleatorização.

  \item[Mensagem do utilizador (Tarefa + Contexto):]
\begin{verbatim}
Generate 15 tasks for:
Subject: Mobilidade Urbana | Level: Beginner
Language: pt-PT
Note: 3 tasks already exist — generate semantically
distinct content.
\end{verbatim}
\end{description}

Se a resposta do LLM contiver JSON malformado ou campos em falta,
o adaptador realiza \textbf{uma única nova tentativa} com o
\textit{prompt} original acrescido de um parágrafo \texttt{schemaReminder}
que recapitula o esquema obrigatório.
Se persistir, o tópico transita para \texttt{FAILED} e um contador
Prometheus \texttt{ai.generation.failures} é incrementado.

\paragraph{INDEXING.}
Os \textit{embeddings} vectoriais de cada tarefa gerada são calculados
e persistidos na coluna \textit{pgvector} de \texttt{task\_templates}
para uso futuro na desduplicação e na pesquisa de similaridade
(explicado na Secção~\ref{sec:pgvector}).

\paragraph{ACTIVE.}
O estado do tópico transita para \texttt{ACTIVE} e um evento SSE
\texttt{complete} é emitido — o portal web apresenta imediatamente o
tópico como disponível para inscrição.

\subsection*{Observação do progresso em tempo real (SSE)}

Após receber o \texttt{201}, o portal web abre uma ligação
\textit{Server-Sent Events}:

\begin{verbatim}
GET /admin/topics/{id}/progress
Accept: text/event-stream
Cache-Control: no-cache
\end{verbatim}

O Nginx é configurado com \texttt{proxy\_buffering off} e
\texttt{X-Accel-Buffering: no} para este \textit{endpoint}, garantindo
que os eventos chegam ao navegador imediatamente sem acumulação no
\textit{proxy}.
O servidor emite eventos no formato W3C SSE~\cite{w3c2021sse}:

\begin{verbatim}
event: phase-change
data: {"phase":"GENERATION"}

event: generation-progress
data: {"completed":7,"total":15}

event: complete
data: {}
\end{verbatim}

O \textit{hook} React \texttt{useTopicProgress} lê o fluxo com a API
\texttt{ReadableStream} (sem \textit{EventSource} de terceiros),
renderizando uma barra de progresso em tempo real.

% ═══════════════════════════════════════════════════════════════════
\section{Pilar II — Personalização: Pesquisa de Similaridade Vectorial}
\label{sec:pgvector}
% ═══════════════════════════════════════════════════════════════════

Esta secção explica a componente mais original do sistema do ponto de vista
da engenharia de dados: a utilização de \textit{embeddings} e pesquisa de
similaridade vectorial para reutilizar progressivamente conteúdo de
remediação entre aprendentes.

\subsection*{O que é um \textit{embedding}}

Um \textit{embedding} é a representação de um texto como um vector de
números reais num espaço de alta dimensão — no caso do \textit{Play4Change},
1024 dimensões produzidas pelo modelo de \textit{embedding} do Mistral.
Textos semanticamente semelhantes (e.g., duas perguntas sobre o mesmo
conceito de mobilidade urbana, formuladas de forma diferente) produzem
vectores próximos nesse espaço.
Textos sobre temas distintos produzem vectores distantes.

\begin{figure}[htbp]
  \centering
  \begin{tikzpicture}[>=Stealth, font=\small]
    % 2D projection of embedding space
    \draw[->] (-0.5,0) -- (7.5,0) node[right]{\small dim$_1$};
    \draw[->] (0,-0.5) -- (0,5.5) node[above]{\small dim$_2$};

    % Cluster A - mobility
    \filldraw[blue!60]  (2.1,3.8) circle(3pt) node[right,font=\tiny]{tarefa mobilidade A};
    \filldraw[blue!40]  (2.5,4.1) circle(3pt) node[right,font=\tiny]{tarefa mobilidade B};
    \filldraw[blue!80]  (1.9,3.5) circle(3pt) node[right,font=\tiny]{dificuldade mobilidade};
    \draw[dashed, blue!40] (2.2,3.8) ellipse (0.8cm and 0.5cm);
    \node[blue!70, font=\tiny] at (2.2,2.95) {cluster mobilidade};

    % Cluster B - energy
    \filldraw[orange!70] (5.5,1.5) circle(3pt) node[right,font=\tiny]{tarefa energia A};
    \filldraw[orange!50] (5.2,1.2) circle(3pt) node[right,font=\tiny]{tarefa energia B};
    \draw[dashed, orange!50] (5.35,1.35) ellipse (0.55cm and 0.45cm);
    \node[orange!70, font=\tiny] at (5.35,0.7) {cluster energia};

    % Cosine angle
    \draw[thick, blue!70,->]  (0,0) -- (2.1,3.8) node[midway,left,font=\tiny]{$\vec{a}$};
    \draw[thick, blue!50,->]  (0,0) -- (2.5,4.1) node[midway,right,font=\tiny]{$\vec{b}$};
    \draw[<->, green!50!black] (0.6,1.1) arc (60:58.5:1.25cm)
      node[right,font=\tiny]{$\theta$ pequeno → sim. elevada};
  \end{tikzpicture}
  \caption{Projeção 2D do espaço de \textit{embeddings}: textos similares formam clusters}
  \label{fig:embedding_space}
\end{figure}

\subsection*{Similaridade do cosseno}

A métrica de proximidade usada é a \textbf{similaridade do cosseno},
definida como:

\[
  \text{sim}(\vec{a}, \vec{b})
    = 1 - \frac{\vec{a} \cdot \vec{b}}{\|\vec{a}\| \cdot \|\vec{b}\|}
    \in [0,\, 1]
\]

onde o valor 1 indica vectores idênticos (ângulo de 0°) e o valor 0
indica vectores ortogonais (ângulo de 90°, sem relação semântica).
A vantagem desta métrica sobre a distância euclidiana é que é
invariante à magnitude dos vectores — o comprimento absoluto do vector
não influencia o resultado, apenas a sua direção no espaço semântico.

O \textit{pgvector} implementa esta operação com o operador binário
\texttt{<=>} que calcula a \emph{distância} do cosseno ($1 - \text{sim}$);
a similaridade obtém-se subtraindo a um:

\begin{verbatim}
SELECT branch_id,
       1 - (se.embedding <=> ?::vector) AS similarity
FROM   struggle_events se
JOIN   adaptive_branches ab ON ab.struggle_event_id = se.id
WHERE  ab.status = 'COMPLETED'
ORDER  BY se.embedding <=> ?::vector
LIMIT  1;
\end{verbatim}

O cast \texttt{?::vector} converte o \textit{array} de \textit{floats}
para o tipo nativo \textit{pgvector}, necessário porque o JDBC não
conhece este tipo de extensão.

\subsection*{O índice IVFFlat: pesquisa aproximada de vizinho mais próximo}

Uma pesquisa exacta por similaridade em $n$ vectores de 1024 dimensões
requer $O(n)$ operações de produto interno — aceitável para centenas de
vectores, mas inviável à escala de milhões.
O \textit{pgvector} implementa o índice \textbf{IVFFlat}
(\textit{Inverted File Flat})~\cite{johnson2019billion} que reduz esta
complexidade ao custo de uma pequena imprecisão (pesquisa
\emph{aproximada}).

\begin{figure}[htbp]
  \centering
  \begin{tikzpicture}[>=Stealth, font=\small]
    % Partition cells
    \foreach \x/\y/\c in {1.5/3.5/blue!15, 4/2/orange!15, 6.5/3.5/green!15,
                           2/1/purple!10, 5/4.5/yellow!20} {
      \filldraw[fill=\c, draw=gray!40] (\x,\y) ellipse (1.1cm and 0.7cm);
    }
    % Centroids
    \foreach \x/\y in {1.5/3.5, 4/2, 6.5/3.5, 2/1, 5/4.5} {
      \filldraw[black] (\x,\y) circle (2.5pt);
    }
    \node[font=\tiny] at (1.5,3.5) [below=4pt] {centróide};

    % Query vector
    \filldraw[red!80] (4.3,2.4) circle(3pt) node[right,font=\tiny]{$\vec{q}$ (nova dificuldade)};

    % Nearest cluster highlighted
    \draw[thick, orange!70, dashed] (4,2) ellipse (1.1cm and 0.7cm);
    \node[font=\tiny, orange!80, align=center] at (4,0.9) {cluster mais próximo\\(único inspeccionado)};

    \draw[->, red!70, thick] (4.3,2.4) -- (4,2);
  \end{tikzpicture}
  \caption{Funcionamento do índice IVFFlat: apenas o cluster mais próximo é inspeccionado}
  \label{fig:ivfflat}
\end{figure}

O IVFFlat divide o espaço vectorial em \texttt{lists} partições durante
a construção do índice (usando \textit{k-means}).
Na pesquisa, apenas as \texttt{probes} partições mais próximas do vector
de consulta são inspeccionadas.
Com \texttt{lists=100} e \texttt{probes=10}, a pesquisa examina apenas
$\approx 10\%$ dos vectores, com uma taxa de revocação tipicamente
superior a 95\%~\cite{johnson2019billion}.

\subsection*{As três estratégias de reutilização}

Quando é detetada uma dificuldade num aprendente, o sistema calcula o
\textit{embedding} da dificuldade e pesquisa o ramo de remediação
mais similar já armazenado.
O resultado da pesquisa determina a estratégia a seguir
(Figura~\ref{fig:reuse_decision}):

\begin{figure}[htbp]
  \centering
  \begin{tikzpicture}[>=Stealth, font=\small,
    box/.style={draw, fill=#1, rounded corners=4pt,
                minimum width=3.2cm, minimum height=0.9cm, align=center},
    dec/.style={draw, diamond, fill=yellow!20,
                minimum width=2.5cm, minimum height=0.9cm, align=center}
  ]
    \node[box=blue!10]    (emb)   at (0, 0)    {Embeder dificuldade\\(1024 dim.)};
    \node[box=blue!10]    (srch)  at (0,-1.8)  {Pesquisa ANN\\pgvector IVFFlat};
    \node[dec]            (sim)   at (0,-3.6)  {similaridade?};

    \node[box=green!18]   (full)  at (-4.5,-5.4){FULL\_REUSE\\{\small sim $>$ 0{,}90}\\{\tiny 0 chamadas ao LLM}};
    \node[box=yellow!22]  (part)  at (0,-5.4)   {PARTIAL\_REUSE\\{\small 0{,}65–0{,}90}\\{\tiny LLM parcial}};
    \node[box=orange!18]  (fresh) at ( 4.5,-5.4){FRESH\_GENERATION\\{\small $<$ 0{,}65}\\{\tiny LLM completo}};

    \node[box=gray!12]    (pers)  at (0,-7.2)   {Persistir ramo +\\embedding};

    \draw[->] (emb)  -- (srch);
    \draw[->] (srch) -- (sim);
    \draw[->] (sim)  -| node[above,font=\tiny]{elevada}  (full);
    \draw[->] (sim)  -- node[right,font=\tiny]{média}    (part);
    \draw[->] (sim)  -| node[above,font=\tiny]{baixa}    (fresh);
    \draw[->] (full)  |- (pers);
    \draw[->] (part)  -- (pers);
    \draw[->] (fresh) |- (pers);
  \end{tikzpicture}
  \caption{Decisão de reutilização com base na similaridade do cosseno}
  \label{fig:reuse_decision}
\end{figure}

\begin{description}
  \item[\texttt{FULL\_REUSE} (similaridade $> 0{,}90$).]
    O ramo de remediação existente é servido diretamente ao aprendente
    sem qualquer chamada ao LLM.
    Este é o \emph{percurso quente}: custo zero de \textit{tokens},
    latência mínima, conteúdo já validado.
    O contador Prometheus
    \texttt{ai.adaptive.reuse\{strategy="full"\}} é incrementado.

  \item[\texttt{PARTIAL\_REUSE} (0{,}65–0{,}90).]
    Uma fração do ramo existente (proporcional à similaridade) é reutilizada;
    as restantes tarefas são geradas de novo e o ramo composto é persistido.
    Este percurso é adequado quando o tema é o mesmo mas o contexto
    específico difere.

  \item[\texttt{FRESH\_GENERATION} ($< 0{,}65$).]
    Não existe nenhum ramo suficientemente similar.
    O \texttt{StruggleAnalysisPrompt} é invocado com o contexto de
    diagnóstico completo e o novo ramo gerado é persistido — enriquecendo
    a biblioteca para utilizadores futuros.
\end{description}

\subsection*{O volante de inércia da personalização}

O mecanismo de reutilização cria um efeito de volante de inércia
(Figura~\ref{fig:flywheel}): cada nova dificuldade diagnosticada
aumenta a biblioteca de ramos disponíveis, tornando as futuras pesquisas
de similaridade mais provável de retornar um resultado de alta qualidade
sem custo de geração.

\begin{figure}[htbp]
  \centering
  \begin{tikzpicture}[>=Stealth, font=\small]
    \def\R{2.8}
    \draw[thick, ->, gray!50] (0:\R) arc (0:355:\R);
    \node[draw, fill=blue!10, rounded corners=4pt, align=center,
          minimum width=2.2cm, minimum height=0.8cm]
          at ( 90:\R) {Aprendente com\\dificuldade};
    \node[draw, fill=orange!12, rounded corners=4pt, align=center,
          minimum width=2.2cm, minimum height=0.8cm]
          at (  0:\R) {Padrão de erro\\classificado};
    \node[draw, fill=purple!10, rounded corners=4pt, align=center,
          minimum width=2.2cm, minimum height=0.8cm]
          at (270:\R) {Ramo gerado ou\\reutilizado};
    \node[draw, fill=green!12, rounded corners=4pt, align=center,
          minimum width=2.2cm, minimum height=0.8cm]
          at (180:\R) {Ramo + \textit{embedding}\\persistidos};
    \node[font=\bfseries\small, align=center] at (0,0)
          {Volante de\\Personalização};
  \end{tikzpicture}
  \caption{O volante de inércia: cada dificuldade enriquece a biblioteca de reutilização}
  \label{fig:flywheel}
\end{figure}

Num sistema a frio (sem dados históricos), toda a geração é inédita.
À medida que a base de utilizadores cresce, a taxa de
\texttt{FULL\_REUSE} aumenta organicamente, reduzindo o custo de
\textit{tokens} e a latência, e melhorando a consistência pedagógica
da remediação.
É o efeito de rede aplicado à personalização pedagógica.

\subsection*{Desduplicação na geração de tarefas}

O mesmo mecanismo vectorial é usado durante a geração de novas tarefas
para evitar conteúdo semântico repetido no mesmo módulo.
Após cada tarefa ser gerada pelo LLM, o adaptador chama
\texttt{PgVectorDeduplicationService.isDuplicate()}, que executa uma
consulta de similaridade contra as tarefas já armazenadas do mesmo tópico.
Uma similaridade superior a 0{,}90 descarta a nova tarefa.
Isto garante diversidade de conteúdo sem necessidade de análise manual
pelo administrador.

% ═══════════════════════════════════════════════════════════════════
\section{Pilar II — Personalização: Deteção de Dificuldade e Remediação}
\label{sec:pilar_remediacao}
% ═══════════════════════════════════════════════════════════════════

A deteção de dificuldade fecha o ciclo da personalização: o sistema
aprende com o comportamento do aprendente em tempo real e adapta
imediatamente a sua jornada de aprendizagem.

\subsection*{Quando é ativada a remediação}

O servidor regista o número de tentativas erradas em
\texttt{TaskAssignment.wrongAttemptCount}.
Quando este valor atinge 2 e o aprendente não tem ainda um ramo
adaptativo atribuído para esta tarefa, o
\texttt{TaskService} invoca
\texttt{HandleStruggleService.triggerAsync()} de forma assíncrona —
a resposta de submissão é devolvida imediatamente ao utilizador sem
aguardar a geração.

\subsection*{Classificação do padrão de erro}

O padrão de erro é inferido a partir da opção selecionada repetidamente
e da proximidade semântica com a opção correta:

\begin{table}[ht]
  \caption{Padrões de erro e directivas do \texttt{StruggleAnalysisPrompt}}
  \label{tab:padroes_erro}
  \vskip 0.5em
  \centering
  \begin{tabular}{lp{8cm}}
    \toprule
    \textbf{Padrão} & \textbf{Directiva de diagnóstico} \\
    \midrule
    \texttt{WRONG\_CONCEPT}        & Re-explicar o conceito central de forma diferente \\
    \texttt{PARTIAL\_UNDERSTANDING}& Focar nos passos intermédios que o aprendente falha \\
    \texttt{KNOWLEDGE\_GAP}        & Preencher o conhecimento pré-requisito antes de regressar \\
    \texttt{UNCLEAR\_INSTRUCTIONS} & Fornecer orientação mais explícita \\
    \texttt{TIME\_PRESSURE}        & Usar tarefas passo-a-passo ao ritmo do aprendente \\
    \texttt{UNKNOWN}               & Andaimamento geral do básico ao intermédio \\
    \bottomrule
  \end{tabular}
\end{table}

\subsection*{Fluxo de remediação na aplicação móvel}

\begin{figure}[htbp]
  \centering
  \begin{tikzpicture}[>=Stealth, font=\small,
    scr/.style={draw, fill=#1, rounded corners=4pt,
                minimum width=3.0cm, minimum height=0.9cm, align=center},
    api/.style={draw, fill=gray!10, rounded corners=3pt,
                minimum width=4.6cm, minimum height=0.7cm, align=center,
                font=\ttfamily\tiny}
  ]
    \node[scr=orange!12] (t1) at (0, 0)   {Tarefa — 2ª tentativa\\errada};
    \node[api]           (s1) at (6.0, 0) {POST /tasks/\{id\}/submit\\struggleTriggered=true};
    \node[scr=red!10]    (t2) at (0,-2.0) {Ecrã Struggle\\{\tiny lista de tarefas adaptativas}};
    \node[api]           (s2) at (6.0,-2.0){GET /struggle/enrollment/\{id\}};
    \node[scr=orange!10] (t3) at (0,-3.8) {Tarefa adaptativa 1};
    \node[api]           (s3) at (6.0,-3.8){POST struggle/.../submit};
    \node[scr=orange!10] (t4) at (0,-5.6) {Tarefa adaptativa N};
    \node[api]           (s4) at (6.0,-5.6){POST .../submit sessionResolved=true};
    \node[scr=green!12]  (t5) at (0,-7.4) {Home\\{\tiny replaceAll}};

    \draw[->] (t1) -- node[above,font=\tiny]{submissão} (s1);
    \draw[->] (s1.west) -- ++(-0.5,0) |- (t2.east);
    \draw[->] (t2) -- node[above,font=\tiny]{carregar} (s2);
    \draw[->] (s2.west) -- ++(-0.5,0) |- (t3.east);
    \draw[->] (t3) -- node[above,font=\tiny]{responder} (s3);
    \draw[->] (s3.west) -- ++(-0.5,0) |- (t4.east);
    \draw[->] (t4) -- node[above,font=\tiny]{responder} (s4);
    \draw[->] (s4.west) -- ++(-0.5,0) |- (t5.east);
  \end{tikzpicture}
  \caption{Sequência de ecrãs e pedidos HTTP durante uma sessão de remediação}
  \label{fig:struggle_sequence}
\end{figure}

Quando a sessão de dificuldade fica resolvida
(\texttt{sessionResolved = true}), o efeito
\texttt{StruggleEffect.NavigateToHome} é emitido.
O \texttt{DefaultRootComponent} usa \texttt{navigation.replaceAll(Config.Home)}
em vez de \texttt{pop()}, substituindo toda a pilha para que o botão de
retrocesso não regresse ao ecrã de dificuldade.

% ═══════════════════════════════════════════════════════════════════
\section{Pilar III — Reconhecimento: Emissão de \textit{Badges}}
\label{sec:pilar_reconhecimento}
% ═══════════════════════════════════════════════════════════════════

O pilar de Reconhecimento transforma conquistas em artefactos verificáveis.
O \texttt{BadgeIssuanceService} é invocado após cada submissão correta e
avalia critérios configuráveis armazenados na tabela \texttt{badge\_criteria}:

\begin{itemize}
  \item \textbf{Conclusão de módulo} — todas as tarefas de um tópico
    concluídas com sucesso.
  \item \textbf{Sequência diária} — \textit{streak} de $N$ dias consecutivos.
  \item \textbf{Limiar de pontuação} — total de pontos acima de um valor
    configurável.
\end{itemize}

Os \textit{badges} emitidos são persistidos na tabela \texttt{user\_badges}
com a data de conquista, o critério cumprido e o identificador do tópico
associado.
O aprendente visualiza os seus \textit{badges} tanto na aplicação móvel
(ecrã de Perfil) como no portal web de administração (vista de detalhe
do utilizador).

% ═══════════════════════════════════════════════════════════════════
\section{Arquitetura do Cliente Móvel: Kotlin Multiplatform}
\label{sec:kmp}
% ═══════════════════════════════════════════════════════════════════

A aplicação móvel partilha a totalidade do código de negócio, rede e
interface entre Android e iOS através de um único projeto Kotlin
Multiplatform.
Apenas os componentes que dependem de APIs nativas sem abstração KMP
existem em conjuntos de fontes específicos por plataforma.

\subsection*{Estrutura de conjuntos de fontes}

\begin{figure}[htbp]
  \centering
  \begin{tikzpicture}[>=Stealth, font=\small,
    common/.style={draw, fill=blue!8, rounded corners=4pt,
                   minimum width=10.5cm, minimum height=0.78cm,
                   align=center, text width=10cm},
    plat/.style={draw, fill=#1, rounded corners=4pt,
                 minimum width=4.8cm, minimum height=0.78cm, align=center}
  ]
    \node[common] at (0, 0.0) {\textbf{commonMain — Interface}
      \ Compose Multiplatform: ecrãs, animações, sistema de design};
    \node[common] at (0,-0.9) {\textbf{commonMain — Apresentação}
      \ Decompose (\texttt{DefaultRootComponent}, \texttt{ChildStack}), efeitos MVI};
    \node[common] at (0,-1.8) {\textbf{commonMain — Domínio}
      \ Repositórios (interface), modelos, \texttt{AppError}, \texttt{NetworkError}};
    \node[common] at (0,-2.7) {\textbf{commonMain — Dados}
      \ \texttt{HttpClientFactory} (Ktor), \texttt{TokenStorage} (interface), DTOs};
    \node[common] at (0,-3.6) {\textbf{commonMain — DI}
      \ Koin: \texttt{AppModule}, módulos por funcionalidade,
      \texttt{platformModule} (expect)};

    \node[plat=green!14] at (-2.6,-5.0)
      {\textbf{androidMain}\\OkHttp, \texttt{EncryptedSharedPreferences}\\
       \texttt{MainActivity}, \texttt{NetworkDetection}};
    \node[plat=gray!16]  at ( 2.6,-5.0)
      {\textbf{iosMain}\\Darwin / NSURLSession, \texttt{KeychainTokenStorage}\\
       \texttt{MainViewController}, \texttt{NWPathMonitor}};

    \foreach \y in {-0.45,-1.35,-2.25,-3.15} {
      \draw[->, gray!50] (5.5,\y) -- (5.5,\y-0.45);
    }
    \draw[->, gray!50] (-2.6,-4.05) -- (-2.6,-4.6);
    \draw[->, gray!50] ( 2.6,-4.05) -- ( 2.6,-4.6);
  \end{tikzpicture}
  \caption{Estrutura KMP: \texttt{commonMain} partilhado e conjuntos de fontes por plataforma}
  \label{fig:kmp_estrutura}
\end{figure}

\subsection*{Navegação com Decompose e padrão MVI}

A navegação é gerida pelo Decompose através de uma
\texttt{StackNavigation<Config>} serializada.
A \texttt{Config} é uma interface selada com oito alternativas
(\texttt{Splash}, \texttt{Login}, \texttt{Home}, \texttt{Task},
\texttt{Profile}, \texttt{About}, \texttt{Explore}, \texttt{Struggle}).
A serialização pelo \texttt{kotlinx.serialization} permite restaurar o
estado da pilha após mudanças de configuração Android sem
\texttt{ViewModel} adicional.

Cada ecrã segue o padrão \textbf{MVI} com três canais:

\begin{figure}[htbp]
  \centering
  \begin{tikzpicture}[>=Stealth, font=\small,
    b/.style={draw, fill=#1, rounded corners=4pt,
              minimum width=3.2cm, minimum height=0.9cm, align=center}
  ]
    \node[b=blue!10]   (view) at (0,    0)  {Ecrã Compose\\(\textit{View})};
    \node[b=green!12]  (comp) at (5.5,  0)  {Componente\\(\textit{Presenter})};
    \node[b=orange!12] (repo) at (5.5,-2.5) {Repositório\\(\textit{Model})};
    \node[b=gray!12]   (state)at (0,   -2.5){\texttt{Value<State>}\\imutável};

    \draw[->, thick] (view)  -- node[above,font=\tiny]{\textit{Events} (onEvent)} (comp);
    \draw[->, thick] (comp)  -- node[right,font=\tiny]{chamadas \texttt{suspend}} (repo);
    \draw[->, thick] (repo)  -- node[below,font=\tiny]{resultado / erro} (comp);
    \draw[->, thick] (comp)  -- node[left, font=\tiny]{\texttt{updateState\{\}}} (state);
    \draw[->, thick] (state) -- node[left, font=\tiny]{\texttt{subscribeAsState()}} (view);
    \draw[->, dashed, red!60] (comp) to[bend left=25]
            node[above,font=\tiny]{\textit{Effects} (navegar)} (view);
  \end{tikzpicture}
  \caption{Fluxo de dados MVI: \textit{Events} → \textit{State} → \textit{Effects}}
  \label{fig:mvi_fluxo}
\end{figure}

O \texttt{BaseView} renderiza um \texttt{ErrorDialog} automaticamente
quando \texttt{state.error != null}, sem lógica adicional nos ecrãs
individuais.
O mapeamento de erros segue a cadeia:
\texttt{Ktor/HTTP} → \texttt{NetworkError} → \texttt{AppError}
→ \texttt{messageKey} localizado — cada camada é agnóstica da camada
abaixo.

% ═══════════════════════════════════════════════════════════════════
\section{Observabilidade: Micrometer, Prometheus e Grafana}
\label{sec:dd_observabilidade}
% ═══════════════════════════════════════════════════════════════════

A observabilidade é a capacidade de inferir o estado interno de um sistema
a partir das suas saídas externas~\cite{humble2018accelerate}.
Num sistema com um subsistema de IA de latência não determinística, a
observabilidade não é um acrescento ao processo de implantação mas um
requisito de concepção.

\subsection*{A pilha de observabilidade}

\begin{figure}[htbp]
  \centering
  \begin{tikzpicture}[>=Stealth, font=\small,
    b/.style={draw, fill=#1, rounded corners=4pt,
              minimum width=3.2cm, minimum height=0.9cm, align=center}
  ]
    \node[b=blue!10]   (app)  at (0, 0)   {Spring Boot\\Micrometer};
    \node[b=orange!12] (prom) at (4.5, 0) {Prometheus\\scraping 5 s};
    \node[b=green!12]  (graf) at (9.0, 0) {Grafana\\Dashboards};
    \node[b=red!10]    (alert)at (4.5,-2.2){Alertmanager};

    \draw[->] (app)  -- node[above,font=\tiny]{/actuator/prometheus} (prom);
    \draw[->] (prom) -- node[above,font=\tiny]{PromQL} (graf);
    \draw[->] (prom) -- node[left, font=\tiny]{alert.rules.yml} (alert);
    \draw[->, dashed] (alert) -- node[right,font=\tiny]{\textit{webhook}} (4.5,-3.4);
  \end{tikzpicture}
  \caption{Pilha de observabilidade: Micrometer → Prometheus → Grafana}
  \label{fig:observabilidade}
\end{figure}

\subsection*{Métricas de domínio}

\begin{table}[ht]
  \caption{Métricas Prometheus personalizadas do \textit{Play4Change}}
  \label{tab:dd_metricas}
  \vskip 0.5em
  \centering
  \begin{tabular}{llp{5.5cm}}
    \toprule
    \textbf{Métrica} & \textbf{Tipo} & \textbf{Descrição} \\
    \midrule
    \texttt{ai.generation.duration}      & Temporizador &
      Duração das chamadas ao LLM por fase \\
    \texttt{ai.generation.failures\_total} & Contador &
      Falhas por tipo (\texttt{schema}, \texttt{adaptive}) \\
    \texttt{ai.adaptive.reuse}           & Contador &
      Distribuição \texttt{full} / \texttt{partial} / \texttt{fresh} \\
    \texttt{ai.dedup.duplicates\_skipped}& Contador &
      Tarefas rejeitadas pela desduplicação \\
    \texttt{struggle\_sessions\_total}   & Contador &
      Novas sessões por \texttt{error\_pattern} \\
    \texttt{struggle\_sessions\_created\_total} & Contador &
      Sessões por tópico \\
    \texttt{topic\_enrollments\_total}   & Contador &
      Inscrições por tópico \\
    \bottomrule
  \end{tabular}
\end{table}

\subsection*{Regras de alerta}

\begin{table}[ht]
  \caption{Regras de alerta Prometheus (8 regras em 4 grupos)}
  \label{tab:alertas}
  \vskip 0.5em
  \centering
  \begin{tabular}{lll}
    \toprule
    \textbf{Alerta} & \textbf{Severidade} & \textbf{Condição} \\
    \midrule
    \texttt{ServerDown}                      & crítico & \texttt{up == 0} por 1 min \\
    \texttt{HighHttpErrorRate}               & aviso   & erros 5xx $> 5\%$ por 2 min \\
    \texttt{HighAiGenerationFailureRate}     & aviso   & $> 0{,}1$ falhas/s por 5 min \\
    \texttt{AiCircuitBreakerOpen}            & crítico & disjuntor aberto (imediato) \\
    \texttt{AiGenerationP95High}             & aviso   & latência p95 $> 90$ s por 5 min \\
    \texttt{JvmHeapUsageHigh}               & aviso   & \textit{heap} $> 85\%$ por 5 min \\
    \texttt{JvmHeapUsageCritical}           & crítico & \textit{heap} $> 95\%$ por 2 min \\
    \texttt{GenerationThreadPoolQueueSaturated} & aviso & fila $\geq 20$ por 2 min \\
    \bottomrule
  \end{tabular}
\end{table}

% ═══════════════════════════════════════════════════════════════════
\section{Modelo de Segurança}
\label{sec:dd_seguranca}
% ═══════════════════════════════════════════════════════════════════

\begin{table}[ht]
  \caption{Mapeamento dos controlos de segurança para o OWASP Top 10 (2021)}
  \label{tab:owasp}
  \vskip 0.5em
  \centering
  \begin{tabular}{lp{9cm}}
    \toprule
    \textbf{Risco} & \textbf{Controlo implementado} \\
    \midrule
    A01 Controlo de acesso quebrado &
      \texttt{@AuthenticationPrincipal} vincula cada recurso ao principal
      autenticado; \texttt{/admin/**} requer papel \texttt{ADMIN} \\
    A02 Falhas criptográficas &
      SHA-256 para \textit{magic link}; HTTPS/TLS obrigatório; sem
      credenciais em texto simples \\
    A03 Injeção &
      Todo o acesso à base de dados via JPA ou \texttt{JdbcTemplate}
      parametrizado; valores pgvector com \texttt{?::vector} \\
    A05 Configuração incorreta &
      Cabeçalhos de segurança Nginx; CORS restrito a origens conhecidas;
      Actuator apenas na porta interna \\
    A07 Falhas de autenticação &
      \textit{Magic link} sem palavra-passe; token de utilização única;
      limitação de taxa \\
    A09 Registo e monitorização &
      Registos estruturados SLF4J + MDC; alertas Prometheus sobre erros \\
    \bottomrule
  \end{tabular}
\end{table}

Adicionalmente, as respostas de erro de domínio não expõem pilhas de
chamadas ou estado interno — o \texttt{@RestControllerAdvice} mapeia
todas as exceções para o esquema \texttt{\{field, reason\}} sem campos
de diagnóstico em produção.

% ═══════════════════════════════════════════════════════════════════
\section{Superfície Completa da API REST}
\label{sec:dd_api_rest}
% ═══════════════════════════════════════════════════════════════════

\begin{table}[ht]
  \caption{Superfície completa da API REST do \textit{Play4Change}}
  \label{tab:api_completa}
  \vskip 0.5em
  \centering
  \footnotesize
  \begin{tabular}{lll}
    \toprule
    \textbf{Método e Caminho} & \textbf{Ator} & \textbf{Finalidade} \\
    \midrule
    \texttt{POST /auth/magic-link}               & Qualquer & Solicitar \textit{magic link} \\
    \texttt{GET  /auth/verify}                   & Qualquer & Redirecionar para \textit{deep link} \\
    \texttt{POST /auth/magic-link/verify}        & Qualquer & Verificar token → par JWT \\
    \texttt{POST /auth/refresh}                  & Qualquer & Renovar \textit{token} de acesso \\
    \texttt{POST /auth/logout}                   & Qualquer & Revogar \textit{token} de atualização \\
    \midrule
    \texttt{POST /admin/topics}                  & Admin. & Criar tópico a partir de URL \\
    \texttt{POST /admin/topics/pdf}              & Admin. & Criar tópico a partir de PDF \\
    \texttt{GET  /admin/topics}                  & Admin. & Listar tópicos (paginado) \\
    \texttt{GET  /admin/topics/\{id\}}           & Admin. & Detalhes e registo de fases \\
    \texttt{GET  /admin/topics/\{id\}/progress}  & Admin. & Fluxo SSE de geração \\
    \texttt{POST /admin/topics/\{id\}/regenerate}& Admin. & Reiniciar \textit{pipeline} \\
    \texttt{GET  /admin/topics/\{id\}/tasks}     & Admin. & Listar modelos de tarefas \\
    \texttt{PUT  /admin/tasks/\{id\}}            & Admin. & Editar modelo de tarefa \\
    \texttt{GET  /admin/topics/\{id\}/struggle-tasks} & Admin. & Listar tarefas adaptativas \\
    \texttt{GET  /admin/topics/\{id\}/prerequisites}& Admin. & Obter pré-requisitos \\
    \texttt{POST /admin/topics/\{id\}/prerequisites}& Admin. & Definir pré-requisitos (DFS) \\
    \texttt{GET  /admin/learning-graph}          & Admin. & DAG completo de pré-requisitos \\
    \texttt{GET  /admin/users}                   & Admin. & Listar utilizadores \\
    \texttt{GET  /admin/users/\{id\}}            & Admin. & Detalhes e estatísticas \\
    \texttt{POST /admin/users/\{id\}/promote}    & Admin. & Promover para administrador \\
    \midrule
    \texttt{POST /user/enroll}                   & Aprendente & Inscrever em tópico \\
    \texttt{GET  /tasks/today}                   & Aprendente & Obter tarefa do dia \\
    \texttt{POST /tasks/\{id\}/submit}           & Aprendente & Submeter resposta \\
    \texttt{GET  /user/roadmap}                  & Aprendente & Mapa de aprendizagem \\
    \texttt{GET  /user/topics}                   & Aprendente & Tópicos disponíveis \\
    \texttt{GET  /struggle/enrollment/\{id\}}    & Aprendente & Sessão de dificuldade \\
    \texttt{POST /struggle/\{s\}/tasks/\{t\}/submit}& Aprendente & Submeter tarefa adaptativa \\
    \texttt{GET  /user/badges}                   & Aprendente & \textit{Badges} conquistados \\
    \texttt{GET  /user/profile}                  & Aprendente & Perfil e pontuação \\
    \texttt{PUT  /user/profile/name}             & Aprendente & Atualizar nome \\
    \texttt{GET  /public/stats}                  & Público & Estatísticas da página inicial \\
    \bottomrule
  \end{tabular}
\end{table}

\paragraph{Referências arquiteturais.}
A arquitectura hexagonal é descrita por Cockburn~\cite{cockburn2005hexagonal}
e aplicada com \textit{Domain-Driven Design}~\cite{evans2003ddd}.
O tratamento funcional de erros usa o tipo \texttt{Either} da
Arrow~\cite{arrow2024}.
O \textit{streaming} SSE segue a especificação W3C~\cite{w3c2021sse}.
A pesquisa vectorial usa a extensão \textit{pgvector}~\cite{pgvector2024}
com indexação IVFFlat~\cite{johnson2019billion}.
A engenharia de \textit{prompts} segue o Guia de Prompting da
Google~\cite{google2024promptguide}.
Os ramos adaptativos fundamentam-se na teoria de andaimamento de
Vygotsky~\cite{vygotsky1978mind} e na aprendizagem por mestria de
Bloom~\cite{bloom1984mastery}.
